\section{Local-to-global deformation results}
\label{app:deformations}\label{sec:global-theorems}

\subsection{Local rings and branch markings}

For the $E_2$ cone the full current ring, including the nilpotent
operator $U$, is \cite{Cremonesi:2015lsa}
\begin{equation}
 \frac{\CC[J_+,J_0,J_-,U]}
 {(J_0^2-J_+J_-,U^2,UJ_+,UJ_0,UJ_-)}.
 \label{eq:e2-full-ring}
\end{equation}
Every weight-zero monomial not containing $U$ reduces to a power
of $J_0$; every nonzero monomial containing $U$ is proportional
to $U$. The Cartan invariants are therefore
\begin{equation}
 \frac{\CC[J_0,U]}{(U^2,UJ_0)}
 \simeq\frac{\CC[\beta,\alpha]}{(\alpha^2,\alpha\beta)}.
 \label{eq:e2-altmann}
\end{equation}
The right side is Altmann's miniversal deformation ring of the
degree-seven cone \cite{Altmann:1997}. Its root marking permits
\begin{equation}
 \alpha=U,\qquad\beta=J_0+cU
 \label{eq:e2-dictionary}
\end{equation}
with compatible nonzero rescalings. The shear preserves the ring
because $U^2=UJ_0=0$. The nilpotent direction vanishes at every
ordinary complex point and is not a second reduced Higgs branch.

For $S_6=\operatorname{Bl}_3\mathbb P^2$, the reduced Higgs
scheme is the union at the origin of the minimal nilpotent orbit
closures of $\mathfrak{sl}_3$ and $\mathfrak{sl}_2$
\cite{Cremonesi:2015lsa}. Their torus-invariant rings are
$\CC[x_1,x_2]$ and $\CC[l]$, giving
\begin{equation}
 \CC[\mathcal H_{E_3}]^{T_\CC}
 =\CC[x_1,x_2]\times_\CC\CC[l]
 =\frac{\CC[x_1,x_2,l]}{(x_1l,x_2l)}.
 \label{eq:e3-invariants}
\end{equation}
This equals the miniversal ring of the degree-six cone
\cite{Altmann:1997}. The two components select the sublattices
\begin{equation}
 A_2=\langle e_1-e_2,e_2-e_3\rangle,\qquad
 A_1=\langle h-e_1-e_2-e_3\rangle
 \subset K_{S_6}^\perp.
 \label{eq:e3-roots}
\end{equation}
To verify the invariant statement and balanced lifts, write a
rank-one traceless matrix as $q_i\widetilde q_j$ with
$z_i=q_i\widetilde q_i$ and $\sum_i z_i=0$. Its torus
invariants are generated by the diagonal entries: every closed
monomial cycle reduces to their product by rank one. Choosing
$|q_i|=|\widetilde q_i|=\sqrt{|z_i|}$ and the appropriate phases
sets the real ADHM and flavor Cartan moment maps to zero.
For the $A_2$ root coordinates take
$(z_1,z_2,z_3)=(x_1,x_2-x_1,-x_2)$; the $A_1$ case uses
two pairs. This supplies the local lifts used in
Proposition~\ref{prop:physical-kernel}.

\subsection{Links and the obstruction map}

Let $\Sigma_p$ be the singularity link and $F_{p,\sigma_p}$ a
Milnor fiber on the chosen branch. Define
\begin{equation}
 V_{p,\sigma_p}=\ker\bigl(H_2(\Sigma_p;\QQ)
                       \longrightarrow H_2(F_{p,\sigma_p};\QQ)\bigr),
 \qquad
 K_\sigma=\ker\bigl(\textstyle\bigoplus_pV_{p,\sigma_p}
                       \longrightarrow H_2(\wideX;\QQ)\bigr).
 \label{eq:milnor-kernel}
\end{equation}
For a node the link space and $V$ are both $\QQ$.
For a del Pezzo cone the Gysin sequence of the unit-circle bundle
of $K_{S_p}$ gives $H_2(\Sigma_p;\QQ)=K_{S_p}^\perp$.
The branch subspaces are the $A_1$ line for $S_7$ and the
$A_1$ line or $A_2$ plane for $S_6$. These are exactly the
cycles killed in the corresponding Milnor fibers
\cite[Proposition~2.1 and Section~2.4]{Wuebben:2026obstruction}.
Thus $K_\sigma\otimes\CC=\ker Q_{X,\sigma}$.

We use the following local comparison in the proof of
Theorem~\ref{thm:first-order-image}. For the chosen local crepant
resolution with exceptional set $E_p$, let
\begin{equation}
 K_p^{\rm cr}=\ker\bigl(H^2_{E_p}(\widehat X_p;\Omega^2)
                       \longrightarrow H^2(\widehat X_p;\Omega^2)\bigr).
\end{equation}
Contraction with the crepant volume form and the boundary map give
compatible identifications
\begin{equation}
 T^1_{X,p}\simeq K_p^{\rm cr}\simeq H_2(\Sigma_p;\CC),
 \label{eq:local-hodge-identification}
\end{equation}
carrying a branch tangent space to $V_{p,\sigma_p}\otimes\CC$.
The local-cohomology maps commute with pushforward of link
classes. On the pentagon the smoothing line is the unique
$A_1$ root line; on the hexagon the reflection and Coxeter
eigenspaces give the $A_1$ and $A_2$ summands. The local
comparison and its volume-form determinant twist are proved in
\cite[Lemmas~2.3--2.4]{Wuebben:2026smoothing}, using
Schlessinger's comparison in the depth-three form
\cite{Schlessinger:1971,Friedman:2022}.

\begin{proof}[Proof of Theorem~\ref{thm:first-order-image}]
For these normal Gorenstein threefolds, the cotangent-complex
and differential deformation groups agree through degree two. The
non-lci locus is finite \cite[Lemma~2.2]{Wuebben:2026smoothing}.
The local-to-global sequence is
\begin{equation}
 0\longrightarrow H^1(X,\Theta_X)\longrightarrow\mathbb T_X^1
 \xrightarrow{\rm loc}\bigoplus_pT^1_{X,p}
 \xrightarrow{\rm ob}H^2(X,\Theta_X).
 \label{eq:local-to-global-ext}
\end{equation}
Excision on $X_{\rm reg}=\wideX\setminus E$ gives
\begin{equation}
 H^1(X_{\rm reg},\Omega^2)\longrightarrow
 \bigoplus_pH^2_{E_p}(\widehat X_p,\Omega^2)
 \longrightarrow H^2(\wideX,\Omega^2).
 \label{eq:global-support-sequence}
\end{equation}
Under \eqref{eq:local-hodge-identification}, the last map followed
by de Rham comparison is the homology pushforward, followed
by Poincar\'e duality. Its source is
$\bigoplus_pH_2(\Sigma_p;\CC)$ and its target is
$H_2(\wideX;\CC)$. A class in $K_\sigma\otimes\CC$ has
zero de Rham image. The $\partial\bar\partial$-lemma for the
projective resolution makes $H^2(\wideX,\Omega^2)\to H^4(\wideX;\CC)$
injective, so it lifts through \eqref{eq:global-support-sequence}.
Contraction with the trivial canonical form identifies the first
group there with $H^1(X_{\rm reg},\Theta)$, and the depth-three
comparison identifies this with $\mathbb T_X^1$.
Naturality gives precisely the localization map in
\eqref{eq:local-to-global-ext}. This proves one inclusion.
Conversely, a global deformation lies in $\ker({\rm ob})$;
the same comparison sends its branch link class to zero in
$H_2(\wideX;\CC)$. This proves the other inclusion.
The admissible-polytope hypotheses ensure that ambient divisors
span the full rational divisor space, so their finite pairing
matrix imposes every continuous constraint
\cite[Proposition~8.3]{Wuebben:2026obstruction}.
\end{proof}

\subsection{Exact smoothing, selected bases and the multiplet change}

\begin{theorem}[mixed smoothing criterion]\label{thm:all-orders-necessity}
Let $X$ be connected, normal and projective with
$\omega_X\cong\mathcal O_X$ and $h^1(X,\mathcal O_X)=0$,
with only nodes and exact analytic anticanonical $S_7,S_6$ cones.
Choose a smooth crepant resolution, small at the nodes and standard
divisorial at the cones, and a reduced local profile $\sigma$.
Then $X$ has a convergent one-parameter smoothing on $\sigma$ if
and only if $K_\sigma$ contains a vector nonzero on every rank-one
summand and with all three coefficients $\mu_i$ nonzero on each
$A_2$ summand, written as $\sum_i\mu_i e_i$ with $\sum_i\mu_i=0$.
Necessity applies to arbitrary analytic smoothing arcs.
\end{theorem}

This is \cite[Theorem~4.7]{Wuebben:2026smoothing} in the degree-six/seven
case. The three $A_2$ coefficients are the nodal smoothing coordinates
on its simultaneous partial resolution. In the root coordinates used
here they are $(H_1,H_2-H_1,-H_2)$, so the plane discriminant is
$H_1H_2(H_1-H_2)=0$. The nodal period argument on the partial model
proves their separate nonvanishing at every order; the scalar-period
argument of \cite{Wuebben:2026obstruction} handles the rank-one germs.
A finite union of proper linear subspaces cannot cover $K_\sigma$, so
checking the individual scalar functionals is sufficient to find one
relation avoiding all discriminants.

For analytic local classifying maps, let $\operatorname{Def}^{\sigma}(X)$
be the scheme-theoretic inverse image of the selected reduced local
branches. The companion integrability theorem states that, if
$K_\sigma\to V_{p,\sigma_p}$ is nonzero at every $E_3$ point, then
\begin{equation}
 \operatorname{Def}^{\sigma}(X)\text{ is smooth},\qquad
 \dim\operatorname{Def}^{\sigma}(X)=h^1(X,\Theta_X)+\dim K_\sigma,
 \label{eq:selected-base-dimension}
\end{equation}
and every tangent on it is realized by a convergent curve
\cite[Theorem~4.5]{Wuebben:2026smoothing}. No extra projection hypothesis
is required at nodes or $E_2$ points. Lifting a relation outside the
discriminants and integrating it proves sufficiency in the preceding
theorem. This does not assert smoothness of unrestricted deformation
schemes, integrability on all zero-projection $E_3$ profiles, or an
identification of a nonlinear classifying map with a linear inclusion.
The geometric theorem also allows degree-five and
$\mathbb P^1\times\mathbb P^1$ cones, but their physical operator
dictionaries are not used in this paper.

For one $E_2$ point there is a direct topological count that
avoids any assumption about the unknown hypermultiplet metric.
Let $n$ be the number of nodes, $k=\dim K_\sigma$ and
$\rho=n+1-k$.

\begin{proposition}\label{thm:forced-spectrum}
If the exact criterion holds and $X_t$ is a smoothing, then
\begin{equation}
 h^{1,1}(\wideX)-h^{1,1}(X_t)=\rho+1,\qquad
 h^{2,1}(X_t)-h^{2,1}(\wideX)=k.
 \label{eq:forced-spectrum-c1}
\end{equation}
\end{proposition}
\begin{proof}
Use the common complement $Y$, exceptional neighborhoods $N_p$
and Milnor fibers $F_p$. The local data are
\begin{equation}
\begin{array}{c|rrrrr}
 &b_2(N_p)&\chi(N_p)&b_2(\Sigma_p)&b_2(F_p)&\chi(F_p)\\\hline
 \text{node}&1&2&1&0&0\\
 C(S_7)&3&5&2&1&1.
\end{array}
\end{equation}
For $S_7$, the Milnor fiber is
$\operatorname{Bl}_{pt}\mathbb P^3\setminus S_7$, and its
link-map kernel is exactly the root line; see the integral
restriction calculation in Appendix~\ref{app:integral-charges}.
The link maps into $N_p$ are injective and $H_1(\Sigma_p;\QQ)=0$,
so $b_2(\wideX)=b_2(Y)+1$. Excision and Lefschetz duality give
$H_3(\wideX,Y;\QQ)=\bigoplus_pH^3(N_p;\QQ)=0$, so
$H_2(Y;\QQ)$ injects into $H_2(\wideX;\QQ)$. The joint link map into
$Y\sqcup\bigsqcup_pF_p$ has kernel $K_\sigma$, hence rank
$n+2-k$. Therefore
\begin{equation}
 b_2(X_t)=b_2(Y)+1-(n+2-k)
         =b_2(\wideX)-\rho-1.
\end{equation}
Euler additivity gives $\chi(X_t)-\chi(\wideX)=-2n-4$.
Using $\chi=2(h^{1,1}-h^{2,1})$ proves the second equality.
\end{proof}
